\section{Introduction}
An observation can be compressed only relative to questions and an
information pattern. When the question is known before compression,
one may retain only its answer. When it arrives later, the same machine
must preserve the possibility of answering several incompatible
questions. With feedback, the retained information also determines which
experiment is performed next. The resulting problem is neither a static
covering problem nor a collection of separate control problems.

This paper studies the marked decision spectrum of a causal experiment.
Its encoder is fixed before the task, its controller and decoder may
know the task but not the unknown parameter, and its persistent state is
charged throughout the execution. The spectrum records the attainable
risk vectors and their dependence on a resource signature. Our principal
result is a controlled common-encoder variational theorem. It constructs
the risk body from the finite instrument itself, characterizes its
feasible continuation partitions, and identifies its lower support
function by an exact occupancy recursion. The maximum over tasks is
retained; it is not replaced by one favourable scalarization.

To describe the result, let $E$ be a finite controlled hidden-state
instrument on a finite horizon, and let $D$ be a finite family of bounded
decision tasks. A row $j\in J$ records a parameter within a minimax task,
or the declared prior of a Bayesian task. Put $M=|J|$. A deterministic
design $\omega$ consists of one common causal encoder together with one
observation-admissible controller and decoder for each task. Write
$R_j(\omega)$ for its row risk, and $b_j$ for the corresponding task's
full-history optimal baseline. These definitions are made precise in
Section~\ref{sec:v9-minimax}; in particular, a minimax task has one
baseline, not a parameterwise oracle.

\begin{theorem}[The common-encoder variational principle]
\label{thm:v9-main}
For prescribed finite register widths ${\bf S}$, let
$\Omega_{\bf S}$ be the deterministic designs. Under an independent
shared seed, averaged before the worst task and parameter are selected,
\begin{align}
 \Theta_{\rm sh}({\bf S};D)
 &=\min_{\lambda\in\Delta(\Omega_{\bf S})}
       \max_{j\in J}\left\{\sum_\omega\lambda_\omega
                                            R_j(\omega)-b_j\right\}
                                                        \label{eq:v9-mainprimal}\\
 &=\max_{w\in\Delta_J}\left\{
           V_0(z_0(w))-w\cdot b\right\}.              \label{eq:v9-maindual}
\end{align}
Here $V_0$ is the finite prescription recursion
\eqref{eq:v9-terminalbellman}--\eqref{eq:v9-bellman}, whose update table
is common to all tasks and parameters. Its support values determine the
whole attainable upper risk body. Deterministic feasible encoders are
exactly the controlled successor-compatible history partitions.

An optimal mixture has at most $M+1$ designs. Removing the shared seed
while retaining its design index in the actual register gives
\[
 \Theta_{\rm pr}((M+1)S;D)
       \le\Theta_{\rm sh}(S;D)\le\Theta_{\rm pr}(S;D).
\]
Private randomization at unchanged width instead has the nonconvex flow
characterization \eqref{eq:v9-privateflow}; it need not give the convex
body in \eqref{eq:v9-mainprimal}.
\end{theorem}

The complete proof is given in Theorems~\ref{thm:v9-minimax} and
\ref{thm:v9-occupancy} and Propositions~\ref{prop:v9-partitions} and
\ref{prop:v9-private}. The underlying convex minimax and separation
arguments are classical. The theorem uses them with a coupled feasible
set: the encoder is common across tasks, each controller is common
across its parameter rows, and a remembered randomization index is an
actual resource when it is not an explicitly permitted shared input.
These constraints are indispensable. In Example~\ref{ex:v9-privategap},
the two-label minimax error of a three-symbol experiment is $1/3$ with a
shared seed and $1/2$ without one. Thus replacing private kernels by an
unpriced convexification would compute a different spectrum.

The same variational principle has two useful specializations. For a
family of Bayesian quadratic tasks it is exactly a maximization over
one task weight vector of the least conditional variance of a weighted
posterior-mean vector, subject to causal compatibility. This strengthens
the single-objective partition formula to the common-encoder maximum
criterion. It applies directly to the accepted-cell/failure protocol
of the A1 manuscript, with a single update on the observed failure
symbol. For exponential criteria it is a duality for moments before
taking logarithms, with a multiplicative occupancy recursion. In a
hidden-phase example it excludes optimizing separately with knowledge
of the phase. These are specializations of one feasible execution body,
not independent assumptions that their controllers can be chosen jointly.

We also construct a least likelihood-saturated predictive coordinate
for every finite instrument. A backward refinement of posterior vectors,
next-report probabilities and successor classes gives its exact minimal
widths. For rational input, this construction and its reconstruction
kernel are computable by finite arithmetic. It makes the finite
sufficiency bridge constructive, while identifying its possible cost:
a one-class future-report quotient can require $2^d$ saturated classes
after $d$ parameter-revealing observations. Statistical saturation is
therefore not a cost-free completion of a predictive quotient.

Finite synthesis yields control of several resources at once.
Theorem~\ref{thm:v9-compiler} compiles an optimal finite design into a
finite program and a dyadic initial lottery. It charges program length,
persistent labels, initialization bits, transient work and bit-serial
running time. Its excess differs from the optimum by at most $k2^{-p}$
for $k\le M+1$ designs and $p$ lottery bits. Synthesis can be exponential
in the horizon; no multiplicative-energy oracle is used in this finite
source theorem. The predictive-tree compiler retained later addresses
a different, infinite symbolic input class under its explicit geometric
and energy-certification assumptions.

A continuous Gaussian example shows that the same causal constraint has
quantitative force beyond finite enumeration. Acquire $d$ noisy scalar
coordinates before an independent coordinate query is disclosed. The
final target is one scalar parameter, but its optimal causal excess is
$S^{-2/d}$, whereas a query-aware static checkpoint has excess $S^{-2}$.
The lower bound uses the common register \emph{before} the query, the
conditional-variance identity, and a density-volume bound on the vector
of posterior means. This is a temporal penalty, not the terminal
dimension argument for an append-only Gaussian vector experiment.

\subsection{Relation to existing theory}
Comparison of experiments by all bounded decisions is the theory of
Blackwell and Le Cam \cite{Blackwell1953,LeCam1964,Torgersen1991}.
The terminal deficiency representation included below is stated with its
compactness hypotheses and the common encoder outside the task
supremum. It does not claim a new unconstrained randomization criterion.
Filtered experiment comparison is itself an established subject
\cite{Norberg2002}. Likewise, predictive equivalence and continuation
congruences have the classical antecedents
\cite{ShaliziCrutchfield2001,Nerode1958}.

Real-time source coding and decentralized control are particularly close
neighbours \cite{Witsenhausen1979,LinderYuksel2014,WoodLinderYuksel2017}.
Nayyar, Mahajan and Teneketzis \cite{NayyarMahajanTeneketzis2013} derive
common-information prescription recursions. In their preprint, Theorem~3
gives the dynamic programme, Corollary~5 treats absent common memory
with finite local memory, and Theorem~4 simplifies its information state.
Our offline occupancy is in this established designer/common-information
tradition. The assertion made here is the coupled minimax risk-body
identity, its controlled partition converse, and the priced passage
between shared and private randomization. An ordinary designer's dynamic
programme, by itself, is not asserted to be new.

The retained quantitative sections use functional and separated-cylinder
quantization, symbolic trees, nonlinear filtering and experiment
approximation \cite{LuschgyPages2004,Roychowdhury2014,KZ2017,Tunstall1967,DKYAC2024}.
Their geometric hypotheses remain local to those theorems. The new
conditional-variance variational principle includes the finite partition
problem and explains the delayed-query Gaussian obstruction. It does
not identify Wasserstein law presentation with decision-complete
statistical deficiency, which retain different output marks.

\subsection{Organization and scope}
Section~\ref{sec:v8-spectrum} fixes the marked spectrum and its transport
and terminal comparison properties. Sections~\ref{sec:v9-minimax} and
\ref{sec:v9-variational} prove the central variational principle.
Sections~\ref{sec:v9-saturation}--\ref{sec:v9-delayed} construct saturation,
compile finite sources and exhibit the Gaussian temporal penalty.
The subsequent sections retain the complete sufficient-statistic,
partition, predictive-tree, measure-presentation, Gaussian and count
arguments. The companion development additionally retains the earlier
mathematical bodies and historical introductions with their original
hypotheses; none is silently replaced by a stronger assertion.

The general marked spectrum includes standard Borel experiments, whereas
the exact controlled minimax theorem proved here has finite horizon,
finite hidden/report/control spaces and finitely many decision rows.
The Gaussian example is established separately by its displayed
continuous argument. In particular, the factor $M+1$ is not a uniform
statement about a growing or infinite task family. A finite-horizon
exponential moment formula is not an LDP, an unbounded-operator theorem,
or a proof of a microscopic phase decomposition. Those further
structures require their own model theorems. The present foundation
supplies a controlled common-encoder converse without importing such
theorems as unnamed premises.
